\section{Tercera línea principal: coevolución de Benchmark, Environment y RL}

\subsection{Sin un benchmark solid, es difícil formar un paradigma unificado}
La discusión sobre benchmark en el panel vale la pena señalarla por separado. Los panelistas no se conforman con «tener varios conjuntos de tareas», sino que insisten en que el benchmark debe ser lo bastante solid como para cubrir distintas capacidades del agent, la retroalimentación del entorno y las dimensiones de evaluación. Solo cuando la amplitud y la dificultad del benchmark son suficientes, la comunidad de investigación puede juzgar qué mejoró realmente el método actual, dónde sigue fallando y qué capacidades son solo un sobreajuste a un rango reducido de tareas.

\begin{importantbox}{Benchmark cumple en el ámbito de Agent un papel más importante que en las tareas de puro texto}
La razón es que Agent involucra más variables:
\begin{itemize}
  \item La entrada ya no es solo el prompt, sino que también incluye el estado del entorno;
  \item La salida ya no es solo texto, sino que también incluye la action sequence;
  \item La evaluación ya no se limita a la respuesta, sino que también observa si el proceso es estable, reproducible y recuperable;
  \item El éxito y el fracaso suelen depender de la environment fidelity.
\end{itemize}
\end{importantbox}

\subsection{Environment es tanto el objeto de entrenamiento como la infraestructura de evaluación}
Varios panelistas, al discutir RL, colocan a environment en un lugar muy alto. El entorno no es un simple telón de fondo, sino la fuente común de reward, el espacio de exploración, tool space y la dificultad del benchmark. Especialmente en tareas como GUI, mobile agent y computer use, si el entorno es estático, falso o no reproducible, tanto el entrenamiento como la evaluación se distorsionan.

\begin{knowledgebox}{Por qué la investigación en Agent se parece cada vez más a la «ingeniería de entornos»}
Porque el scaling of acting no basta y es necesario seguir con el scaling of environment. En la discusión se mencionó explícitamente:
\begin{itemize}
  \item Se necesita más infra, más entornos de tareas y más datos;
  \item Se necesita un tool space y un action space más grandes;
  \item Se necesitan entornos reales, teléfonos/computadoras en la nube o entornos aislados reproducibles;
  \item Se necesita un world model o un environment model para reducir el costo de construcción.
\end{itemize}
\end{knowledgebox}

\subsection{Self-evolving se parece más a una metodología que a una tarea única}
Sobre self-evolving, los panelistas señalaron un punto que se pasa por alto con facilidad: a diferencia de reasoning, que a simple vista corresponde a una capacidad concreta, self-evolving se parece más a un conjunto de metodologías de automejora en torno a distintas tareas. Por eso ámbitos como code, math, GUI y visión difícilmente comparten de inmediato un paradigma completamente unificado. En el corto plazo, la ruta más realista es formar primero un consenso local en cada subámbito, para luego avanzar poco a poco hacia un método unificado más amplio.

\begin{warningbox}{No esperes que todas las tareas de Agent compartan de inmediato un mismo recipe de entrenamiento}
El environment de texto, código, GUI y visión difiere demasiado: la densidad de recompensa, el espacio de acción y los modos de fallo son distintos. Forzar una fórmula unificada tiende a borrar detalles clave y termina en una solución de compromiso que no es lo bastante buena para nadie.
\end{warningbox}

\subsection{Resumen de la sección}
Si la investigación en Agent quiere avanzar hacia un paradigma unificado, no puede eludir benchmark y environment. Ambos determinan tanto si RL es entrenable como si la evaluación es confiable. Más que «buscar otro prompt mejor», esta línea principal se parece más a construir una infraestructura de investigación capaz de iterar de manera continua las capacidades de Agent.

